\documentclass[11pt,a4paper]{article}
\usepackage[utf8]{inputenc}
\usepackage[T1]{fontenc}
\usepackage{amsmath,amssymb,amsthm}
\usepackage{geometry}
\usepackage{titlesec}
\usepackage{fancyhdr}
\usepackage{graphicx}
\usepackage{array}
\usepackage{booktabs}
\usepackage{longtable}

\geometry{margin=1in}

\pagestyle{fancy}
\fancyhf{}
\fancyhead[C]{\small Cayley, Collected Mathematical Papers, Vol.~XI}
\fancyfoot[C]{\thepage}
\renewcommand{\headrulewidth}{0.4pt}

\theoremstyle{plain}
\newtheorem{art}{Art.}[section]

\begin{document}

\begin{center}
\Large\bfseries 745.\\[6pt]
\Large\bfseries ON THE SCHWARZIAN DERIVATIVE\\
\Large\bfseries AND THE POLYHEDRAL FUNCTIONS.\\[12pt]
\normalsize
[From the \textit{Philosophical Transactions}, vol.~CLXXIII (1879),\\
pp.~611--659.]\\[12pt]
\small
[Received November 13, ---Read December 19, 1879.]
\end{center}

\normalsize


The last-mentioned differential equation is considered by Klein and Brioschi:
the solutions in 13 cases, or such of them as had not been given by Schwarz,
were obtained by Brioschi: and those of the remaining 3 cases, subject to
a correction in one of them, were afterwards obtained by Klein.

The first part of the present Memoir relates, say to the foregoing equation
\[
\{s, x\} = (a, b, c \frown \frown) \left(\frac{1}{x-a}, \frac{1}{x-b}, \frac{1}{x-c}\right),
\]
although the other form in $[x, z]$ may equally well be regarded as the
fundamental form. We consider in the theory:

\textbf{A.} The Derivative $\{s, x\}$, meaning as above explained.

\textbf{B.} Quadric functions of any three or more inverts $\frac{1}{x-a}$.

\textbf{C.} Rational and integral functions $P$, $Q$, $R$ having a sum $=0$,
and which are such that $QR' - Q'R$, $RP' - R'P$, $PQ' - P'Q$, contains
only the factors of $P$, $Q$, $R$.

\textbf{D.} The differential equation of the third order.

\textbf{E.} The Schwarzian theory in regard to conformable figures and the
corresponding values of the imaginary variables $s$ and $x$.

\textbf{F.} Connexion with the differential equation for the hypergeometric
series.

The second part of the Memoir relates to the Polyhedral Functions.
The paragraphs of the whole Memoir are numbered consecutively.

\section*{PART I. The Derivative $\{s, x\}$.}

\begin{art}
Art.~Nos.~1 to~7.
\end{art}

\begin{art}
The derivative $\{s, x\}$ may be transformed in regard to either or both of
the variables. Suppose, first, that $s$ is a function of the new variable~$S$,
(hence also $S$ is a function of~$x$): using subscript numbers to denote
differentiations in regard to~$S$, and the accents as before for
differentiations in regard to~$x$, we have
\[
\frac{s'}{S'} = \frac{ds}{dS}, \quad
\frac{s''}{S'^2} - \frac{s'S''}{S'^3} = \frac{d^2s}{dS^2},
\]
whence, differentiating the logarithms,
\[
\frac{s''}{s'} - \frac{S''}{S'} = \frac{1}{\frac{ds}{dS}} \frac{d^2s}{dS^2} \cdot S',
\]
and consequently
\[
\frac{s'''}{s'} - \frac{s''^2}{s'^2} - \frac{S'''}{S'} + \frac{S''^2}{S'^2}
= \frac{d}{dx} \left(\frac{1}{\frac{ds}{dS}} \frac{d^2s}{dS^2}\right) S'^2
+ \frac{1}{\frac{ds}{dS}} \frac{d^2s}{dS^2} \cdot S'';
\]
that is,
\[
\{S, x\} = \{s, x\} - \left(\frac{dS}{dx}\right)^2 \{s, S\}.
\]
\end{art}


\begin{art}
Again differentiating, we have the required formula.

In a very similar manner, taking $x$ a function of~$X$, it is shown that
\[
\{s, X\} = \{s, x\} - \left(\frac{dx}{dX}\right)^2 \{x, X\}.
\]
\end{art}

\begin{art}
If in this formula we write $S$ for~$s$, and substitute the resulting value
of $\{S, x\}$ in the former formula, we have
\[
\{S, X\} = \{s, x\} - \left(\frac{dS}{dx}\right)^2 \{s, S\}
- \left(\frac{dx}{dX}\right)^2 \{x, X\},
\]
which is the formula for the change of both variables.

It, in fact, includes the other two: viz. writing $X = x$, or $S = s$, and
observing that $\{s, s\} = 0 = \{x, x\}$, we have the other two formulae.
\end{art}

\begin{art}
By putting in the first formula $X = x$, we obtain a formula for the
interchange of the variables.
\end{art}

\begin{art}
Writing $S = \frac{as+b}{cs+d}$, and using for a moment the accents to denote
differentiation in regard to~$s$, we have
\[
S' = \frac{ad-bc}{(cs+d)^2}, \quad
\frac{S''}{S'} = -\frac{2c}{cs+d},
\]
and thence
\[
\frac{S'''}{S'} - \frac{3}{2}\left(\frac{S''}{S'}\right)^2 = 0;
\]
that is, $\{S, s\} = 0$, whence also $\{s, S\} = 0$.
\end{art}


\begin{art}
Hence in the first formula $\{S, x\} = \{s, x\}$, that is,
\[
\left\{\frac{as+b}{cs+d},\ x\right\} = \{s, x\};
\]
viz. we may, in the derivative $\{s, x\}$, write for~$s$ any homographic
function $(as+b)/(cs+d)$ of~$s$.
\end{art}

\begin{art}
Again, if $X = \frac{\alpha x + \beta}{\gamma x + \delta}$, then from the
second formula
\[
\{s, X\} = \{s, x\} - \frac{(\alpha\delta - \beta\gamma)^2}{(\gamma x + \delta)^4}
\left\{x,\ \frac{\alpha x + \beta}{\gamma x + \delta}\right\},
\]
that is,
\[
\left\{s,\ \frac{\alpha x + \beta}{\gamma x + \delta}\right\}
= \{s, x\} - \frac{(\alpha\delta - \beta\gamma)^2}{(\gamma x + \delta)^4}
\left\{x,\ \frac{\alpha x + \beta}{\gamma x + \delta}\right\},
\]
and here, changing $s$ into $(aS+b)/(cS+d)$, we have finally
\[
\left\{\frac{as+b}{cs+d},\ \frac{\alpha x + \beta}{\gamma x + \delta}\right\}
= \frac{(\gamma x + \delta)^4}{(\alpha\delta - \beta\gamma)^2} \{S, x\},
\]
which is the formula for the homographic transformation of the two
variables $s$,~$x$.
\end{art}

\begin{art}
Let $s$ be a given function of~$x$, the equation $\{S, x\} = \{s, x\}$ is a
differential equation of the third order in~$S$, and by what precedes, its
general integral is
\[
S = \frac{as+b}{cs+d}.
\]
The direct process is as follows: we have a first integral
\[
\frac{S''}{S'} = \frac{s''}{s'} - \frac{2cs'}{cs+d};
\]
a second integral $\log S' = \log s' - 2\log(cs+d) + \text{const.}$, that is,
\[
S' = \frac{A \cdot s'}{(cs+d)^2};
\]
and thence a final integral $S = B - \frac{A}{c(cs+d)}$, which is equivalent
to the foregoing value of~$S$.
\end{art}

\section*{The Quadric Function of three or more Inverts. Art.~Nos.~8 to~15.}

\begin{art}
We consider a quadric function of any number of inverts
$\frac{1}{x-a}, \frac{1}{x-\alpha}, \ldots$, all of them different: it is
assumed that the constant term is $=0$, and also that the sum of the
coefficients of the linear terms is $=0$. We have therefore square terms
$\frac{a}{(x-a)^2}$, product terms $\frac{h}{(x-a)(x-\beta)}$, and linear
terms $\frac{A}{x-a}$, where the sum of the coefficients $A$ is $=0$.

Any product term $\frac{h}{(x-a)(x-\beta)}$ is expressible in the form of a
difference of two linear terms, $\frac{h}{a-\beta}\left(\frac{1}{x-a} - \frac{1}{x-\beta}\right)$,
and (the coefficients of these being equal), after it is thus expressed,
the sum of the coefficients of the linear terms is still~$=0$.

The function is thus always expressible in the form
\[
\frac{a}{(x-a)^2} + \frac{b}{(x-\beta)^2} + \cdots
+ \frac{A}{x-a} + \frac{B}{x-\beta} + \cdots,
\]
where the sum $A + B + \ldots$ is $=0$: this may be called the \textit{reduced form}.
\end{art}

\begin{art}
Observe that any particular invert may disappear altogether from the
reduced form: this will be the case if $a = 0$, that is, if the original
form contains no term in $\frac{1}{(x-a)^2}$, and if also $A = 0$.

An invert thus disappearing from the reduced form is said to be
\textit{non-essential}: and the inverts which do not disappear are said to be
\textit{essential}. The original form contains in appearance the non-essential
inverts, but it is really a quadric function of the essential inverts only.
\end{art}

\begin{art}
Imagine the original function expressed as a rational function of~$x$, say
\[
\frac{U}{V^2},
\]
where $V = (x-a)(x-\beta)\cdots$ is the product of the linear factors
(corresponding to the several inverts), each taken with the index~$1$; and
$U$ is a function of the same degree as~$V^2$.
\end{art}

\begin{art}
In particular, if the number of essential inverts is $=3$, then the quadric
function is of the form
\[
(a, b, c, f, g, h \,\raisebox{-.5ex}{$\widetilde{\phantom{aa}}$}\, x-a, x-\beta, x-\gamma)^2,
\]
which contains one superfluous constant, and equivalent functions differ
only by a multiple of
\[
\frac{1}{(x-\beta)(x-\gamma)} + \frac{1}{(x-\gamma)(x-\alpha)} + \frac{1}{(x-\alpha)(x-\beta)} = 0.
\]
\end{art}

\begin{art}
A quadric function such that the degree of the numerator is less by~4
than that of the denominator is said to be ``curtate.''

The conditions, in order that the function
\[
(a, b, c, f, g, h \,\raisebox{-.5ex}{$\widetilde{\phantom{aa}}$}\, x-a, x-\beta, x-\gamma)^2
\]
may be curtate, are easily found to be
\begin{align*}
a + b + c + 2f + 2g + 2h &= 0, \\
\alpha(h + b + f) + \beta(g + f + c) + \gamma(a + h + g) &= 0;
\end{align*}
and by reason of the superfluous constant we are at liberty to assume a
third condition: the three conditions may be taken to be
\[
a + h + g, \quad h + b + f, \quad g + f + c \quad \text{each} = 0;
\]
and these being satisfied, the quadric function is curtate.
\end{art}

\begin{art}
The function
\[
\frac{a}{x-\alpha} + \frac{b}{x-\beta} + \frac{c}{x-\gamma} + \cdots,
\]
where the coefficients $a$, $b$, $c$,~\ldots\ satisfy the relation
$a + b + c + \ldots = -2$, is diaphoric, and therefore curtate.

In fact, forming the sum,
\[
\text{coeff.\ } \frac{1}{(x-\alpha)^2}
+ (\text{coeff.\ } \frac{1}{(x-\alpha)(x-\beta)} + \cdots),
\]
this is $-a - \frac{1}{2}a - \frac{1}{2}b - \frac{1}{2}c - \ldots = -\frac{1}{2}(2 + a + b + c + \ldots)$,
which is $=0$; and similarly the other conditions are satisfied.
\end{art}

\begin{art}
The function regarded as a function of the inverts
\[
\frac{1}{x-a}, \frac{1}{x-a_1}, \ldots; \frac{1}{x-b}, \frac{1}{x-b_1}, \ldots; \frac{1}{x-c}, \frac{1}{x-c_1}, \ldots,
\]
where $a + a_1 + \ldots = b + b_1 + \ldots = c + c_1 + \ldots = h$ suppose,
is diaphoric, and therefore curtate.
\end{art}


\section*{Rational and integral functions $P$, $Q$, $R$ having a sum $=0$.
Art.~Nos.~16 to~20.}

\begin{art}
The functions $P$, $Q$, $R$ are functions of a variable~$z$; for the present
purpose they need only be rational and integral functions of~$z$, but they
are really functions connected with the theory of the inverts as will
appear. They are each of them of the same order~$S$; any two functions
not originally of the same order, can thus be made to be of the same
order; or by taking account of the root $z = \infty$, we may in the original
case regard them as being of the same order, and it is convenient so to
regard them: say they are taken to be of the same order~$S$.

And there is clearly no loss of generality in taking the three functions
to be prime to each other; for any common factor of two of them would
divide the third, and might therefore be struck out.
\end{art}

\begin{art}
We may therefore write
\[
P = F \prod (z-l)^p, \quad Q = G \prod (z-m)^q, \quad R = H \prod (z-n)^r,
\]
where $(z-l)^p$ is taken to denote the distinct simple or multiple factors
of~$P$, and the like as regards $Q$ and~$R$; the factors $z-l$, $z-m$, $z-n$
are thus all of them different.

And we have
\[
S = \Sigma p = \Sigma q = \Sigma n.
\]
\end{art}

\begin{art}
It is at once seen that $QR' - Q'R$ is of the degree $2S-2$, and moreover
that it contains the factors $\prod (z-l)^{p-1}$, $\prod (z-m)^{q-1}$,
$\prod (z-n)^{r-1}$; hence it contains the factor
\[
\prod (z-l)^{p-1} (z-m)^{q-1} (z-n)^{r-1}.
\]

Suppose the number of distinct indices~$p$ is $\sigma_1$, that of distinct
indices~$q$ is~$\sigma_2$, and that of distinct indices~$r$ is~$\sigma_3$;
then the degree of the factor is $= 3S - \sigma_1 - \sigma_2 - \sigma_3$;
and if this be $= 2S - 2$, then $\Theta$ can have no other variable factor:
viz. if the numbers $\sigma_1$, $\sigma_2$, $\sigma_3$ of the distinct indices $p$, $q$, $r$
respectively are such that $\sigma_1 + \sigma_2 + \sigma_3 = S + 2$,
a relation which is henceforth taken to be satisfied, then we have
\[
\Theta = K \prod (z-l)^{p-1} (z-m)^{q-1} (z-n)^{r-1}.
\]
\end{art}

\begin{art}
As already in effect remarked, the conclusion extends to the other functions
$RP' - R'P$ and $PQ' - P'Q$; and we thus see that the condition in order
that each of these functions may contain only the factors of $P$, $Q$, $R$ is
\[
\sigma_1 + \sigma_2 + \sigma_3 = S + 2.
\]
\end{art}

\begin{art}
The case most closely connected with the theory of inverts is when each
of the indices~$p$, each of the indices~$q$, and each of the indices~$r$ is
$=1$; the condition then is $\sigma_1 + \sigma_2 + \sigma_3 = S + 2$, where
$S = \sigma_1 = \sigma_2 = \sigma_3$, that is, $3S = S + 2$, or $S = 1$.
But this gives only the case $P : Q : R = 1 : -1 : 0$.
\end{art}

\section*{The Differential Equations $\{x, z\}$ and $\{s, x\}$.
Art.~Nos.~21 to~45.}

\begin{art}
In reference to what follows, it is convenient to put $P = \lambda P_0$,
$P_\lambda = \lambda^{-1} P$, where $P_0$ is written for $\prod (z-l)^{p-1}$, the
G.C.M. of $P$ and~$P'$; and $\lambda$ is consequently $= F$ multiplied by the
product $\prod (z-l)$ of the several factors taken each with the index
unity; and so for $Q$ and~$R$:
viz. we write
\[
P,\ Q,\ R = \lambda P_0,\ \mu Q_0,\ \nu R_0, \quad
P_\lambda,\ Q_\mu,\ R_\nu = \lambda^{-1}P,\ \mu^{-1}Q,\ \nu^{-1}R,
\]
and the foregoing value of $\Theta$ then is $\Theta = K P_0 Q_0 R_0$.

We come now to the investigation of the leading theorem.

Take $a$, $b$, $c$ arbitrary, $f, g, h = b-c,\ c-a,\ a-b$; $P$, $Q$, $R$
functions of~$z$ as above; and write
\[
f(x-a) : g(x-b) : h(x-c) = P : Q : R,
\]
equations, which are consistent with each other and determine~$x$ as a
rational function of~$z$.

Using, as before, the accent to denote differentiation in regard to~$z$,
and taking the coefficients $(a, b, c)$ arbitrary, it is required to find the
value of
\[
\{x, z\} + x'^2 (a, b, c \,\raisebox{-.5ex}{$\widetilde{\phantom{aa}}$}\, x-a, x-b, x-c)^2.
\]
\end{art}

\begin{art}[Calculation of the first term $\{x, z\}$]
We have $x = $ a function $\frac{\alpha P + \beta R}{\gamma P + \delta R}$, and thence
$\{x, z\} = \{s, z\}$ where $s = P/R$ for a moment; then
\[
\frac{P}{R} : \frac{Q}{R} : 1 = \frac{P_0 Q_0 Z}{Z^2 R_0} : \frac{P_0 Q_0}{Z \cdot R_0} : \frac{P_0 Q_0}{Z^2 R_0},
\]
that is,
\[
\frac{P}{R} = \frac{P_0 Q_0 Z}{Z^2 R_0}, \quad
\frac{Q}{R} = \frac{P_0 Q_0}{Z \cdot R_0}, \quad
1 = \frac{P_0 Q_0}{Z^2 R_0}.
\]

Substituting the values
\[
P_0 = \prod (z-l)^{p-1}, \quad Q_0 = \prod (z-m)^{q-1}, \quad
R_0 = \prod (z-n)^{r-1}, \quad Z = \prod (z-n),
\]
we have
\[
\{x, z\} = \frac{P_0''}{P_0} - \frac{Q_0''}{Q_0} + 4\frac{Z'}{Z}
\left(\frac{P_0'}{P_0} - \frac{Q_0'}{Q_0} + \frac{Z'}{Z}\right)
- 2\left(\frac{P_0'}{P_0} - \frac{Q_0'}{Q_0} + \frac{Z'}{Z}\right)^2
\]
\[
- \frac{P_0'}{P_0} \cdot \frac{Q_0'}{Q_0}
+ \left(\frac{P_0'}{P_0} - \frac{Q_0'}{Q_0}\right) \frac{Z'}{Z}
+ 2\left(\frac{Z'}{Z}\right)^2
\]
\end{art}


\begin{art}
where it is to be observed that
\[
\Sigma(p-1) + \Sigma(q-1) - \Sigma(r+1)
= S - \sigma_1 + S - \sigma_2 - (S + \sigma_3) = S - \sigma_1 - \sigma_2 - \sigma_3 = -2;
\]
consequently the function is diaphoric, and therefore curtate.

It is to be remarked that the function, although presenting itself in a
form unsymmetric in regard to the factors of $P$ and~$Q$, and of~$R$,
is really symmetric as regards the three sets of factors; this is obvious
\textit{a priori}, and it will be presently verified.
\end{art}

\begin{art}[Calculation of the second term]
For the calculation of the second term
\[
x'^2 (a, b, c \,\raisebox{-.5ex}{$\widetilde{\phantom{aa}}$}\, x-a, x-b, x-c)^2,
\]
we have
\[
f(x-a),\ g(x-b),\ h(x-c) = P,\ Q,\ R.
\]
\end{art}

\begin{art}
We have to bring the function on the right-hand side into the reduced form
\[
\frac{a}{(z-l)^2} + \frac{A}{z-l},
\]
for the purpose of getting rid of the non-essential inverts (if any).

We write
\[
\frac{P''}{P} = \Sigma \frac{(p-1)}{(z-l)^2} + \Sigma \frac{(p-1)}{z-l} \cdot \frac{P_0'}{P_0} + \cdots
\]
viz. $z-l$ here denotes any particular factor, and $z-l_1$ represents any
other factor of the same set; and so in other like cases.
\end{art}

\begin{art}
The whole coefficient of $\frac{1}{(z-l)^2}$ is
\[
-(p-1) - \frac{1}{2}(p-1)^2 + \frac{1}{2}p(p+1) = \frac{1}{2}(1-p^2) + ap = (ap);
\]
an expression which, regarded as a function of $a$ and~$p$, is represented
by $(ap)$: the parentheses are used only to avoid ambiguity, and are
omitted when there is no ambiguity.
\end{art}

\begin{art}
The whole coefficient of $\frac{1}{z-l}$, when $z-l$ is a multiple factor
of~$P$, thus is
\[
(ap) \cdot \frac{1}{z-l} + \text{terms in other inverts},
\]
a form which is now symmetrical in regard to the inverts.
\end{art}

\begin{art}
The value just obtained must be equal to
\[
\frac{(ap)}{(z-l)^2} + \frac{A}{z-l} + \cdots + x'^2 (a, b, c \,\raisebox{-.5ex}{$\widetilde{\phantom{aa}}$}\, x-a, x-b, x-c)^2;
\]
viz. comparing the two forms and reducing, they will be identical if only
\[
\frac{(ap)}{z-l} + \frac{(bq)}{z-m} + \frac{(cr)}{z-n} = 0,
\]
and it can be shown that the function inside the $\{\}$ is in fact $=0$.
\end{art}

\begin{art}
We have, as before, $\Sigma \frac{p}{z-l} = \Sigma \frac{P'}{P}$; or writing each of
these quantities $= \phi$, the equation to be verified is
\[
\frac{(ap)}{z-l} + \frac{(bq)}{z-m} + \frac{(cr)}{z-n} = 0.
\]
\end{art}

\begin{art}
But from the equation $\Theta = PQ' - P'Q = KP_0 Q_0 R_0$, we find
$XY' - X'Y = KR_0$; and then, differentiating,
$ZF_0' + Z'F_0 - Z_0'F - Z_0F' = K'R_0 + KR_0'$: writing in these equations
$z = l$, they become
\[
-X_0 F = KR_0, \quad X'Y_0 - X_0'Y - X_0 Y' = KR_0',
\]
so that, dividing the second by the first,
\[
\frac{X'Y_0 - X_0'Y - X_0 Y'}{X_0 F} = \frac{R_0'}{R_0},
\]
or, recollecting that $X_0 = pX'$ and $Y = qY_0$, we have
\[
\frac{X'}{X_0} - \frac{Y'}{Y_0} = \frac{R_0'}{R_0},
\]
that is,
\[
\frac{p}{z-l} - \frac{q}{z-m} = \frac{R_0'}{R_0},
\]
the required relation.
\end{art}

\begin{art}
The result is that, $z-l$ being a multiple factor of~$P$, the coefficient
of $\frac{1}{(z-l)^2}$ is $(ap) = \frac{1}{2}(1-p^2) + ap$, and the coefficient
of $\frac{1}{z-l}$ is a sum of terms each containing the factor~$(ap)$.
\end{art}

\begin{art}
The formulae for the coefficients of $\frac{1}{(z-m)^2}$, $\frac{1}{z-m}$, and
$\frac{1}{(z-n)^2}$, $\frac{1}{z-n}$ give at once, by a mere change of letters,
those for the coefficients of $\frac{1}{(z-n)^2}$, $\frac{1}{z-n}$; and the
function in question,
\[
\{x, z\} + x'^2 (a, b, c \,\raisebox{-.5ex}{$\widetilde{\phantom{aa}}$}\, x-a, x-b, x-c)^2,
\]
is now obtained in the required form
\[
\frac{(ap)}{(z-l)^2} + \frac{(bq)}{(z-m)^2} + \frac{(cr)}{(z-n)^2}
+ \frac{A}{z-l} + \frac{B}{z-m} + \frac{C}{z-n},
\]
where $(ap)$ denotes $\frac{1}{2}(1-p^2) + ap^2$, and the like for $(bq)$ and~$(cr)$;
and where, $z-l$ being a multiple factor of~$P$, the coefficient~$A$ contains
the factor~$(ap)$; and similarly for $B$ and~$C$.
\end{art}


\begin{art}
The reasoning applies directly to lines 2, 3, 4, 5 of the $PQR$-Table:
and with a slight variation to line~1; viz. here the factors of $R$
($= -1 + z^n$) are all simple factors, but in virtue of $c = \infty$ and
$a = b$, the corresponding inverts disappear, and, the other inverts also
disappearing, the value of the function is $=0$.

Hence lines 1, 2, 3, 4, 5 of the $PQR$-Table give each of them a result
$=0$, for the values of $(a, b, c)$ appearing by the table itself, and shown
explicitly in the corresponding line of the Annex.

Thus line~3 shows that the function
\[
\{x, z\} + x'^2 (3, -1, -1 \,\raisebox{-.5ex}{$\widetilde{\phantom{aa}}$}\, x-1, x+1, x-\infty)^2
\]
is identically $=0$; and so for the other lines.
\end{art}

\begin{art}
It is hardly necessary to remark that an expression
\[
(a_1, b_1, c_1, \ldots \,\raisebox{-.5ex}{$\widetilde{\phantom{aa}}$}\, x-a_1, x-b_1, x-c_1, \ldots)^2
\]
in fact denotes
\[
\frac{a_1}{(x-a_1)^2} + \frac{b_1}{(x-b_1)^2} + \frac{c_1}{(x-c_1)^2}
+ \frac{2f_1}{(x-a_1)(x-b_1)} + \frac{2g_1}{(x-a_1)(x-c_1)}
+ \frac{2h_1}{(x-b_1)(x-c_1)} + \cdots.
\]
The particular form of the $z$-inverts is immaterial; we could by a
general homographic transformation of the variable change these into any
other set of inverts; and the foregoing is therefore a solution of
\[
\{s, x\} = (a, b, c \,\raisebox{-.5ex}{$\widetilde{\phantom{aa}}$}\, x-a, x-b, x-c)^2,
\]
a differential equation of the third order.

This is the Differential Equation $[a, b, c \ldots]$.
\end{art}

\begin{art}
From the Roman lines, if we assume
\[
f(x-a) : g(x-b) : h(x-c) = \mathfrak{P} : \mathfrak{Q} : \mathfrak{R},
\]
where $\mathfrak{P}$, $\mathfrak{Q}$, $\mathfrak{R}$ are functions of~$z$, not the
same functions that $P$, $Q$, $R$ are of~$s$, since they belong to a
different line of the Table: we have, as before,
\[
\{x, z\} + x'^2 (a, b, c \,\raisebox{-.5ex}{$\widetilde{\phantom{aa}}$}\, x-a, x-b, x-c)^2 = 0.
\]
\end{art}

\begin{art}
We may combine any such result with a properly selected result of the
preceding system, the two results being such that $(a, b, c)$ have the
same values: combining in this manner the several results, we obtain
altogether a system of 16 results.
\end{art}

\begin{art}
It thus appears that there are in all 16 sets of values of $(a, b, c)$,
for which the equation is solved, viz. the 16 sets of values are shown in
the right-hand column of the Annex. For greater clearness I exhibit the
integral equations as follows:

\begin{center}
\begin{tabular}{|l|l|l|l|}
\hline
& Functions of $x$ & & Functions of $s$ \\
\hline
(1) & $f(x-a) : g(x-b) : h(x-c) = P : Q : R$ & Polygon & I \\
(2) & Double Pyramid & II & $x^2 : -(x^2+1)^2 : (x^2-1)^2$ \\
(3) & Tetrahedron & III & $x^3 : -(x+1)^3 : (x-1)^3$ \\
(4) & Cube and Octahedron & V & $(x^4+1)^2 : -(x^4-1)^2 : 4x^4$ \\
(5) & Dodecahedron and Icosahedron & VII & \\
\hline
\end{tabular}
\end{center}
\end{art}

\begin{art}
Consider, in general, a solution of this form, $x = F(s)$ a rational
function of~$s$: then $s$ is an irrational function of~$x$, and if
$s_1$, $s_2$ are any two of its values,
\[
\{s_1, x\} = R(x), \quad \{s_2, x\} = R(x);
\]
that is, $\{s_1, x\} = \{s_2, x\}$, and therefore (ante, No.~7)
\[
s_2 = \frac{as_1 + b}{cs_1 + d};
\]
then $s = f(s_1) = f\left(\frac{as_1+b}{cs_1+d}\right) = F(s_2)$: viz. $F(s)$ is a
rational function of~$s$, transformable into itself by the transformation
$s$ into $\frac{as+b}{cs+d}$; and it is moreover clear that between any two
roots~$s$ whatever of the equation $x = F(s)$ there exists a homographic
relation of the form in question.
\end{art}


\section*{The Schwarzian Theory. Art.~Nos.~46 to~70.}

\begin{art}
I consider the foregoing equation $\{s, x\} = R(x)$, where $R(x)$ is a
rational function, and is now taken to be a real function of~$x$: we may
assume $s' = \rho' e^{i\theta}$, where the accents denote differentiation in regard
to~$x$, and where $\rho'$, $\theta$, and therefore also $\theta'$, are real functions
of~$x$.

We have
\[
\frac{s''}{s'} = \frac{\rho''}{\rho'} + i\theta',
\]
and thence
\[
\{s, x\} = \{\rho, x\} + \{\theta, x\} + i\left(\theta''' - \theta' \frac{\rho''}{\rho'}\right).
\]

Putting this $= R(x)$, and assuming that $x$ is real, we have
\[
\{\rho, x\} + \{\theta, x\} = R(x), \quad \theta''' - \theta' \frac{\rho''}{\rho'} = 0.
\]
\end{art}

\begin{art}
In the $s$-plane: since $\{s, x\} = R(x)$ is a differential equation of the
third order, the integral contains three arbitrary constants, and we may
imagine these so determined that to the values $x = a, b, c$ shall correspond
the values $s = A, B, C$ respectively.

If there is not on the $x$-line any critical point, as the point~$x$ moves
continuously along this line the point~$s$ will move continuously along a
circle.
\end{art}

\begin{art}
If $A$ be the value of~$s$ corresponding to $h = 0$, then $s = hA$, and we find
\[
\frac{s - A}{s - A_1} = \frac{A - A_1}{A - A_2} \cdot \frac{h - h_1}{h - h_2} \cdot e^{i\phi},
\]
viz. reducing $\frac{s-A}{s-A_1}$ to its principal term~$h^\mu$, and then writing
$ds$, $dx$ for $s-A$, and $h (= x-a)$ respectively, we have
\[
ds = K (dx)^\mu,
\]
or $x = a$ is a critical point with the exponent~$\mu$; and similarly $x = b$
and $x = c$ are critical points with the exponents $\lambda$ and~$\nu$ respectively.
\end{art}

\begin{art}
Hence in the equation
\[
\{s, x\} = (a, b, c \,\raisebox{-.5ex}{$\widetilde{\phantom{aa}}$}\, x-a, x-b, x-c)^2,
\]
the exponents $\mu$, $\lambda$, $\nu$ at the critical points $a$, $b$, $c$ are connected
with the coefficients $(a, b, c)$ by the relations
\[
a = \frac{1}{2}(1 - \mu^2), \quad b = \frac{1}{2}(1 - \lambda^2), \quad c = \frac{1}{2}(1 - \nu^2).
\]
\end{art}

\begin{art}
As an instance, take the double pyramid form: the integral equation is
\[
f(x-a) : g(x-b) : h(x-c) = 4s^n : -(s^n-1)^2 : (s^n+1)^2,
\]
or say
\[
\frac{(c-a)(x-b)}{(a-b)(x-c)} = \frac{4s^n}{(s^n+1)^2},
\]
or if, for greater simplicity, we assume $a, b, c = 1, 0, \infty$, this is
\[
x = \frac{4s^n}{(s^n+1)^2},
\]
or say
\[
\frac{1}{\sqrt{x}} = \frac{s^n+1}{2s^{n/2}}, \quad \text{that is, } s^n = \frac{1 \pm \sqrt{1-x}}{1 \mp \sqrt{1-x}},
\]
a solution of the differential equation
\[
\{s, x\} = \frac{1}{2}\frac{1-\frac{1}{4}n^2}{x^2} + \frac{1}{2}\frac{1-\frac{1}{4}n^2}{(x-1)^2} + \frac{\frac{1}{4}n^2 + \frac{1}{4}n^2 - 1}{x(x-1)}.
\]

In particular, if $n = 3$, we have $s^3 = \left(\frac{1+\sqrt{1-x}}{1-\sqrt{1-x}}\right)^2$.
\end{art}

\begin{art}
Take $x$ real; then, if $x$ is positive and less than~1, $s^n$ is real and
positive, and we have for~$s$ the infinite half-lines at the inclinations
$0^\circ$, $120^\circ$, $240^\circ$; while if $x$ is positive and greater than~1,
$s^{n/2}$ is real and negative, and we have the infinite half-lines at the
inclinations $60^\circ$, $180^\circ$, $300^\circ$.

If $x$ is real and negative, then $\frac{1+ki}{1-ki} = \cos\phi + i\sin\phi$: whence
$s$ is of the same form, or the locus of the point~$s$ is a circle.
\end{art}

\begin{art}
We have
\[
\frac{b-c}{(x-a)^2} \cdot \frac{c-a}{(x-b)^2} \cdot \frac{a-b}{(x-c)^2} = FGH \cdot e^{i(\theta_a + \theta_b + \theta_c - 2\pi)},
\]
with the like values for $G$ and~$H$. Hence the right-hand side of
\[
\frac{c-a}{(x-b)(x-c)} \cdot \frac{a-b}{(x-c)(x-a)} \cdot \frac{b-c}{(x-a)(x-b)} = FGH
\]
has the values
\[
\frac{a}{x-a}, \frac{b}{x-b}, \frac{c}{x-c}.
\]
\end{art}

\begin{art}
Considering now the left-hand side of the equation, we have
\[
\{s, x\} = \frac{1}{(x-a)(x-b)(x-c)} \left(\frac{a}{x-a} + \frac{b}{x-b} + \frac{c}{x-c}\right),
\]
substituting for~$x$ its value $= \alpha + \beta i + \gamma e^{i\phi}$, this becomes
\[
\{s, x\} = -\frac{1}{2\rho^2} \left(\{s, \xi\} - \frac{1}{2}\right),
\]
that is,
\[
\{s, x\} = -\frac{1}{2\rho^2} \left(\{s, \xi\} - \frac{1}{2}\right).
\]
\end{art}

\begin{art}
Again $P$, $Q$ given functions of~$z$, and $v$ a function of~$z$ determined
by the equation
\[
\frac{d^2v}{dz^2} + Pv = 0,
\]
and assume $y = uv$. Substituting this value of~$y$ in the first equation,
we obtain for~$v$ an equation of the second order (the coefficients of which
contain~$u$), and we may make this identical with the second equation;
viz. comparing the coefficients of the two equations, we thus have two
equations each containing~$u$; and by eliminating~$u$ we obtain a
differential equation of the third order for the determination of~$z$ as a
function of~$x$.
\end{art}

\begin{art}
We obtain as the required equation for the determination of~$z$ as a
function of~$x$:
\[
\{z, x\} + \left(\frac{dz}{dx}\right)^2 (a_1, b_1, c_1 \,\raisebox{-.5ex}{$\widetilde{\phantom{aa}}$}\, z-\infty, z-0, z-1)^2 = 0.
\]

The process does not give the value of~$u$, but this can be found without
difficulty, viz.
\[
\frac{du}{dx} = \frac{1}{2} \left(\frac{z''}{z'} - \frac{x''}{x'}\right).
\]
\end{art}

\begin{art}
Jacobi's differential equation of the third order for the transformed
modulus~$\lambda$, \textit{Fund.~Nova}, p.~78, [Ges.\ Werke, t.~i, p.~132], is
\[
3(\lambda'^2 \lambda''' - \lambda''^3) - 2\lambda \lambda' (\lambda' \lambda''' - \lambda'' \lambda''') + \lambda'^4 = 0.
\]
\end{art}

\begin{art}
And putting in the formula $x = \infty$, $= -(1-a)$, we have
\[
\frac{d^2y}{dx^2} + \left(\frac{1-a-b-c}{x} + \frac{1-a'+b-c}{x-1}\right) \frac{dy}{dx}
+ \left(\frac{-a}{x^2} + \frac{b}{x(x-1)} + \frac{c}{(x-1)^2}\right) y = 0,
\]
with a like formula for $\frac{d^2y}{dx^2} + 2P \frac{dy}{dx} - 4Qy = 0$.

We then have
\[
y = uv, \quad v = x^{-a_1} (1-x)^{-c_1} z^{a'} (1-z)^{b'+c'+d'+e'} \ldots,
\]
and the differential equation of the third order for the determination
of~$z$ is
\[
\{z, x\} + \left(\frac{dz}{dx}\right)^2 (a_1, b_1, c_1 \,\raisebox{-.5ex}{$\widetilde{\phantom{aa}}$}\, z-\infty, z-0, z-1)^2 = 0,
\]
where $a_1$, $b_1$, $c_1$ are functions of $a$, $b$, $c$,~\ldots
determined as above.
\end{art}

\begin{art}
Consider a spherical surface and upon it any number of points: take at
pleasure any point as South Pole, this determines the plane of the
equator; and the stereographic projection of any point is the intersection
with the plane of the equator of the line joining the point with the South
Pole.

To fix the ideas take the radius of the sphere as unity: let the axes
of $x$ and~$y$ be drawn in the plane of the equator in longitudes $0^\circ$
and $90^\circ$ respectively.
\end{art}


\section*{PART II. The Polyhedral Functions. Art.~Nos.~71 to~117.}

\begin{art}
The values of $x + iy$ for the mid-points of the sides are
\[
\cos\frac{\pi}{n} + i\sin\frac{\pi}{n}, \quad
\cos\frac{3\pi}{n} + i\sin\frac{3\pi}{n}, \quad\ldots,\quad
\cos\frac{(2n-1)\pi}{n} + i\sin\frac{(2n-1)\pi}{n}.
\]
The corresponding function of~$s$ is $s^n + 1$.

The North and South Poles, which form with the $n$~points a double
pyramid of $n+2$ summits, correspond to the values $s = 0$ and $s = \infty$.
We have thus $s(1-s^n)(s^n-1)$ as the function corresponding to the
double pyramid.
\end{art}

\begin{art}
(3). Considering for a moment the tetrahedron as a figure inscribed in a
cube, with the summits at alternate corners of the cube; if the cube is
placed with its faces parallel to the coordinate planes, and with edge
$= \sqrt{2}$, then the coordinates of the 8~corners are $(\pm 1, \pm 1, \pm 1)$,
and those of the 4~tetrahedral summits are $(1,1,1)$, $(1,-1,-1)$,
$(-1,1,-1)$, $(-1,-1,1)$.
\end{art}

\begin{art}
(4). The octahedron is placed with two of its summits as poles, and the
other four summits in the equator at longitudes $0^\circ$, $90^\circ$, $180^\circ$,
$270^\circ$ respectively: the values of~$s$ are, as in the last case, $0$, $\infty$,
$+1$, $\pm i$, and the function is $s(s^4-1)$.

The function for the centres of the faces, or summits of the cube, is
$s^8 + 14s^4 + 1$.

The function for the mid-points of the sides of the octahedron or of the
cube is $s^{12} - 33s^8 - 33s^4 - 1$.
\end{art}

\begin{art}
(5). The Icosahedron is placed with two opposite summits as poles, and
the other 10~summits on two parallel circles of 5~summits each.
\end{art}

\begin{art}
Since $f_4$ is the same function as $f_5$, we have of course $f_4$, $h_4$ and
$t_4$ themselves covariants of~$f_5$: but it is convenient to separate the
two systems.
\end{art}

\begin{art}
It is to be observed that $f_5$ is a quartic function having its
quadrinvariant $(f) = 0$; but independently of this, that is, qua quartic
function, it has only the covariants $h_5$ and~$t_5$ (the Hessian and the
cubicovariant respectively), viz. every other covariant is a rational and
integral function of $f_5$, $h_5$, $t_5$.
\end{art}

\section*{Covariantive Formulae. Art.~Nos.~81 to~84.}

\begin{art}
The various covariantive formulae will be given with their proper
numerical coefficients.

\textbf{Tetrahedron function.} $f$, $h$, $t$ stand for the before-mentioned
values $f_3$, $h_3$, $t_3$ ($P, Q, R = h^2, -12i\sqrt{3} \cdot t, -f^3$).

For $f_3$: $(a, b, c, d, e) = 1, 0, -\frac{1}{2}\sqrt{2}, 0, 1$.

\begin{align*}
i(f,f)^1 &= -96i\sqrt{3} \cdot h, \\
i(f,h)^1 &= 96i\sqrt{3} \cdot f, \\
i(f,t)^1 &= -32f\cdot h, \\
(f,h) &= 32i\sqrt{3} \cdot t, \\
(f,f)^2 &= 576f = 576 \cdot (-\frac{1}{8}), \\
(h,t) &= 4f\cdot h.
\end{align*}
\end{art}

\begin{art}
\textbf{Dodecahedron function.} $f$, $h$, $t$ stand for the before-mentioned
values $f_5$, $h_5$, $t_5$ ($P, Q, R = h^2, -t^2, -1728f^5$).

For $f_5$: $(a, b, c, d, e, f, g, h, i, j, k, l, m) = (0, \frac{1}{2}, 0, 0, 0, 0, \frac{5}{8}, 0, 0, 0, 0, -\frac{5}{8}, 0)$.

\begin{align*}
i(f,f)^1 &= -121h, \\
i(f,f)^2 &= 0, \\
i(f,f)^3 &= (924) = (720) = \frac{1}{4}f^2, \\
i(f,h)^1 &= 0, \\
i(f,h)^2 &= 0, \\
i(f,h)^3 &= (924)(720) = \frac{1}{4}f\cdot h, \\
(f,h) &= -20t, \\
i(h,h)^1 &= 173280f\cdot h, \\
(f,t) &= -30h^2, \\
i(h,t)^1 &= 9082800f^2\cdot h, \\
(h,t) &= -86400f^2, \\
h^2 - t &= -1728f^5 = 0.
\end{align*}
\end{art}

\begin{art}
We have
\[
t = (s^4 + f^4)(1 + 2\sqrt{-1} \cdot s^2 f^2),
\]
where the product of the last two factors is $s^8 + (k^2 + k'^2)s^4 - 1$.

We have
\[
k = \sqrt{\frac{1}{2}}(80\sqrt{5} - 176), \quad k' = \sqrt{\frac{1}{2}}(5\sqrt{5} - 11),
\]
and consequently $k^2 - k'^2 = 11$; or the function is
\[
s(1-s^4)(s^8 + 11s^4 - 1).
\]
\end{art}

\begin{art}
Similarly, writing for shortness $k = \tan\gamma$, $k' = \tan\gamma'$, where
\[
\cos^2\gamma = \frac{3+\sqrt{5}}{8}, \quad \sin^2\gamma = \frac{5-2\sqrt{5}}{15}, \quad
\cos^2\gamma' = \frac{3-\sqrt{5}}{8}, \quad \sin^2\gamma' = \frac{10+2\sqrt{5}}{15},
\]
and therefore $\frac{k^2}{k'^2} = \frac{5-2\sqrt{5}}{10+2\sqrt{5}}$, $\sin^2\gamma = \frac{5-2\sqrt{5}}{15}$, as before, then the
value is
\[
s(s^4 + k^2)(s^4 + k'^2).
\]
\end{art}

\section*{Invariantive property of the Stereographic Projection.
Art.~Nos.~87 to~93.}

\begin{art}
The before-mentioned theorem that the functions derived from two
different stereographic projections of the same point are linear
transformations one of the other, may be thus stated:

Considering on the surface of a sphere, two fixed points A and~B;
and determining the position of a point~$C$, first in regard to~A by its
distance~$d$ and azimuth~$\phi$, and next in regard to~B by its distance~$d_1$
and azimuth~$\phi_1$: the theorem is that $e^{i\phi_1} \tan\frac{1}{2}d_1$ is a homographic
function of $e^{i\phi} \tan\frac{1}{2}d$.
\end{art}

\begin{art}
But it is interesting to give the proof with rectangular coordinates.
Taking $(X, Y, Z)$, $(X_1, Y_1, Z_1)$ for the coordinates, referred to two
different sets of axes, of a point on the spherical surface: also $x$, $y$,
$x_1$, $y_1$ for the coordinates of the corresponding stereographic
projections, we have
\[
(X_1, Y_1, Z_1) = (\alpha, \beta, \gamma)(X, Y, Z),
\]
where $(\alpha, \beta, \gamma)$, $(\alpha', \beta', \gamma')$, $(\alpha'', \beta'', \gamma'')$ are the direction-cosines of
the new axes in regard to the old ones.
\end{art}

\begin{art}
viz. attending only to the former of these, we have $\frac{x_1}{y_1}$ is a homographic
function of $\frac{x}{y}$, which is the before-mentioned theorem.
\end{art}

\begin{art}
Supposing that the transformation from $(X, Y, Z)$ to $(X_1, Y_1, Z_1)$ is
made by a rotation, the coordinates of which are $\lambda$, $\mu$, $\nu$, that is,
if $f$, $g$, $h$ are the inclinations of the resultant axis to the axes of $x$,
$y$, $z$ respectively, and $\theta$ the angle of rotation, putting
\[
\lambda, \mu, \nu = \tan\tfrac{1}{2}\theta \cos f, \quad \tan\tfrac{1}{2}\theta \cos g, \quad \tan\tfrac{1}{2}\theta \cos h,
\]
then
\[
\frac{x_1}{y_1} = \frac{(1+\lambda^2-\mu^2-\nu^2)x + 2(\lambda\mu-\nu)y + \ldots}{2(\lambda\mu+\nu)x + (1-\lambda^2+\mu^2-\nu^2)y + \ldots}.
\]
\end{art}

\begin{art}
$(As + B) \div (Cs + D)$.

The group of 60 and the group of 24 include each of them as part of
itself the group of 12: it is further to be remarked that the group of 12
may be regarded as that of the positive substitutions upon four letters
$a, b, c, d$, the group of 24 as that of all the substitutions upon the four
letters, and the group of 60 as that of the positive substitutions upon
five letters $a, b, c, d, e$.
\end{art}

\begin{art}
I call to mind that a group of functional symbols $1$, $\alpha$, $\beta$, \ldots\ is such
that the product (or compound) of any two of them is a symbol of the
group; the symbols may be representable as rational (or irrational)
functions of a variable symbol, but in the theory of groups this is not
necessary.
\end{art}


\begin{art}
4 points~A and 4 points~B. The fifteen circles intersect by fives in the
pairs of opposite points~A, by threes in the pairs of opposite points~B,
and by twos in the pairs of opposite points~$\Theta$; the mutual inclinations
of successive circles at the points A, B,~$\Theta$ being $=36^\circ$, $60^\circ$ and
$90^\circ$ respectively.

The whole number $15 \cdot 14 = 210$, of the intersections of the circles
two and two together is thus made up of the 12~points~A each counting
10~times, the 20~points~B each counting 3~times, and the 30~points~$\Theta$
each counting once; $210 = 120 + 60 + 30$.
\end{art}

\begin{art}
The angular magnitudes which present themselves are all obtained from
the dodecahedral pentagon, as shown in the annexed figure, in which the
angle subtended by a side at the centre is $= 72^\circ$, and the angle between
adjacent sides is $= 108^\circ$.
\end{art}

\begin{art}
And from the triangle, where the two sides and the included angle are
$g$, $g$ and $(120^\circ - 2\phi = 2\xi)$, we find $\widehat{ae} = 60^\circ$.
\end{art}

\begin{art}
We thus arrive at the following Table:

\begin{center}
\begin{tabular}{|l|l|l|l|l|l|l|}
\hline
& $\sin$ & $\cos$ & & & & \\
\hline
A$\frac{1}{2}\alpha$ & & & & & & \\
B$\frac{1}{2}\beta$ & & & & & & \\
AB & & & & & & \\
$\Theta$B & & & & & & \\
B$\Theta$B & & & & & & \\
BBB & & & & & & \\
$\Theta$$\Theta$ & & & & & & \\
\hline
\end{tabular}
\end{center}

where as above
\begin{align*}
\alpha &= 31^\circ 43', & \beta &= 20^\circ 55', & \gamma &= 37^\circ 22', \\
\chi &= 70^\circ 32', & \theta &= 54^\circ 44', & \delta &= 37^\circ 46', \\
\phi &= 22^\circ 14', & \omega &= 63^\circ 26', & 2\beta &= 41^\circ 50', \\
2\gamma &= 74^\circ 44', & \alpha - \beta &= 10^\circ 48', & \text{etc.}
\end{align*}

and
\[
\alpha + \beta + \gamma = 90^\circ, \quad \chi + 2\theta = 180^\circ, \quad \omega + \phi = 60^\circ.
\]
\end{art}

\begin{art}
We now construct three figures of the points A, B,~$\Theta$; viz. these are
stereographic projections, each showing the Northern hemisphere projected
on the plane of the equator by lines drawn to the South Pole: hence, for
any pair of opposite points not on the equator, only the point in the
Northern hemisphere is shown.
\end{art}

\begin{art}
Fig.~1. Pole at~$\Theta$.

\begin{center}
\begin{tabular}{|l|l|l|}
\hline
& N.P.D.'s & Longitudes \\
\hline
2A & $\alpha = 31^\circ 43'$ & $0^\circ$, $180^\circ$ \\
2A & $90^\circ - \alpha = 58^\circ 17'$ & $90^\circ$, $270^\circ$ \\
1A & $90^\circ - (0^\circ, 180^\circ) + \alpha = 31^\circ 43'$ & \\
2A & $90^\circ + \alpha = 121^\circ 43'$ & $90^\circ$, $270^\circ$ \\
2A & $180^\circ - \alpha = 148^\circ 17'$ & $0^\circ$, $180^\circ$ \\
\hline
$\frac{1}{2}$B & $\beta = 20^\circ 55'$ & $90^\circ$, $270^\circ$ \\
1B & $\theta = 54^\circ 44'$ & $45^\circ$, $135^\circ$, $225^\circ$, $315^\circ$ \\
2B & $90^\circ - \beta = 69^\circ 5'$ & $0^\circ$, $180^\circ$ \\
& $(90^\circ, 270^\circ) + \beta = 20^\circ 55'$ & $45^\circ$ \\
2B & $90^\circ + \beta = 110^\circ 55'$ & $0^\circ$, $180^\circ$ \\
4B & $180^\circ - \theta = 125^\circ 16'$ & $45^\circ$, $135^\circ$, $225^\circ$, $315^\circ$ \\
2B & $180^\circ - \beta = 159^\circ 5'$ & $90^\circ$, $270^\circ$ \\
\hline
$\frac{1}{2}\Theta$ & $0^\circ$ & --- \\
4$\Theta$ & $36^\circ$ & $(90^\circ, 270^\circ) + \alpha = 31^\circ 43'$ \\
\hline
\end{tabular}
\end{center}
\end{art}

\begin{art}
Fig.~2. Pole at~A.

\begin{center}
\begin{tabular}{|l|l|l|}
\hline
& N.P.D.'s & Longitudes \\
\hline
A & --- & --- \\
\hline
$\frac{1}{2}$A & $2\alpha = 63^\circ 26'$ & $0^\circ$, $72^\circ$, $144^\circ$, $216^\circ$, $288^\circ$ \\
$\frac{1}{2}$A & $180^\circ - 2\alpha = 116^\circ 34'$ & $36^\circ$, $108^\circ$, $180^\circ$, $252^\circ$, $324^\circ$ \\
A & $180^\circ$ & --- \\
\hline
$\frac{1}{2}$B & $\gamma = 37^\circ 22'$ & $36^\circ$, $108^\circ$, $180^\circ$, $252^\circ$, $324^\circ$ \\
$\frac{1}{2}$B & $90^\circ - \alpha + \beta = 79^\circ 12'$ & $36^\circ$, $108^\circ$, $180^\circ$, $252^\circ$, $324^\circ$ \\
$\frac{1}{2}$B & $90^\circ + \alpha - \beta = 100^\circ 48'$ & $72^\circ$, $144^\circ$, $216^\circ$, $288^\circ$ \\
$\frac{1}{2}$B & $180^\circ - \gamma = 142^\circ 38'$ & $72^\circ$, $144^\circ$, $216^\circ$, $288^\circ$ \\
\hline
$\frac{1}{2}\Theta$ & $\alpha = 31^\circ 43'$ & $72^\circ$, $144^\circ$, $216^\circ$, $288^\circ$ \\
$\frac{1}{2}\Theta$ & $90^\circ - \alpha = 58^\circ 17'$ & $36^\circ$, $108^\circ$, $180^\circ$, $252^\circ$, $324^\circ$ \\
$\frac{1}{2}\Theta$ & $90^\circ$ & $(36^\circ, 108^\circ, 180^\circ, 252^\circ, 324^\circ) + 18^\circ$ \\
$\frac{1}{2}\Theta$ & $90^\circ + \alpha = 121^\circ 43'$ & $72^\circ$, $144^\circ$, $216^\circ$, $288^\circ$ \\
$\frac{1}{2}\Theta$ & $180^\circ - \alpha = 144^\circ 17'$ & $36^\circ$, $108^\circ$, $180^\circ$, $252^\circ$, $324^\circ$ \\
\hline
\end{tabular}
\end{center}
\end{art}

\begin{art}
Fig.~3. Pole at~B.

\begin{center}
\begin{tabular}{|l|l|l|}
\hline
& N.P.D.'s & Longitudes \\
\hline
$\frac{1}{2}$A & $\gamma = 37^\circ 22'$ & $30^\circ$, $150^\circ$, $270^\circ$ \\
$\frac{1}{2}$A & $90^\circ - \alpha + \beta = 79^\circ 12'$ & $90^\circ$, $210^\circ$, $330^\circ$ \\
$\frac{1}{2}$A & $90^\circ + \alpha - \beta = 100^\circ 48'$ & $30^\circ$, $150^\circ$, $270^\circ$ \\
$\frac{1}{2}$A & $180^\circ - \gamma = 142^\circ 38'$ & $90^\circ$, $210^\circ$, $330^\circ$ \\
\hline
B & --- & --- \\
\hline
$\frac{1}{2}$B & $2\beta = 41^\circ 50'$ & $90^\circ$, $210^\circ$, $330^\circ$ \\
$\frac{1}{2}$B & $\chi = 70^\circ 32'$ & $(30^\circ, 150^\circ, 270^\circ) + \delta = 37^\circ 46'$ \\
$\frac{1}{2}$B & $180^\circ - \chi = 109^\circ 28'$ & $(90^\circ, 210^\circ, 330^\circ) + \delta = 37^\circ 46'$ \\
$\frac{1}{2}$B & $180^\circ - 2\beta = 138^\circ 10'$ & $30^\circ$, $150^\circ$, $270^\circ$ \\
\hline
$\frac{1}{2}\Theta$ & $\beta = 20^\circ 55'$ & $90^\circ$, $210^\circ$, $330^\circ$ \\
$\frac{1}{2}\Theta$ & $\theta = 54^\circ 44'$ & $(90^\circ, 210^\circ, 330^\circ) + \phi = 22^\circ 14'$ \\
$\frac{1}{2}\Theta$ & $90^\circ - \beta = 69^\circ 5'$ & $30^\circ$, $150^\circ$, $270^\circ$ \\
$\frac{1}{2}\Theta$ & $90^\circ$ & $60^\circ$, $120^\circ$, $180^\circ$, $240^\circ$, $300^\circ$ \\
$\frac{1}{2}\Theta$ & $90^\circ + \beta = 110^\circ 55'$ & $90^\circ$, $210^\circ$, $330^\circ$ \\
$\frac{1}{2}\Theta$ & $180^\circ - \theta = 125^\circ 16'$ & $(30^\circ, 150^\circ, 270^\circ) + \phi = 22^\circ 14'$ \\
$\frac{1}{2}\Theta$ & $180^\circ - \beta = 159^\circ 5'$ & $30^\circ$, $150^\circ$, $270^\circ$ \\
\hline
\end{tabular}
\end{center}
\end{art}

\section*{The groups of homographic transformations, resumed.
Art.~Nos.~104 to~117.}

\begin{art}
The axes of rotation for the dodecahedron and the icosahedron are 15~axes
each through a pair of opposite points~$\Theta$, 6~axes each through a pair
of opposite points~A, and 10~axes each through a pair of opposite
points~B; or say 15~$\Theta$-axes, 10~$B$-axes and 6~$A$-axes: the corresponding
angles of rotation are $180^\circ$, $72^\circ$ and $120^\circ$; so that (excluding in
each case the original position or that of a rotation~0) we have in respect
of each $\Theta$-axis 1~position, in respect of each $A$-axis 4~positions, and in
respect of each $B$-axis 2~positions; in all, including the original position,
\[
1 + 15 + (6 \times 4) + (10 \times 2) = 60
\]
positions, that is, a group of 60~rotations.
\end{art}

\begin{art}
To find, in any one of the three forms, the substitution corresponding to
given values of $(\lambda, \mu, \nu)$, that is, to a rotation through an angle~$\theta$
about an axis through the point $(f, g, h)$.
\end{art}

\begin{art}
\textbf{Homographic Transformations. The group of 60.}

Pole at~$\Theta$.

\begin{center}
\begin{longtable}{|c|c|c|c|c|}
\hline
No. & $(As+B) \div (Cs+D)$ & & & \\
\hline
1 & $s$ & 1 & & \\
2 & $-s$ & $-1$ & 1 & \\
3 & $\frac{1}{s}$ & 1 & & \\
4 & $\frac{1}{-s}$ & $-1$ & & \\
5 & $2$ & $-3+\sqrt{5}+i(-1+\sqrt{5})$ & $-3+\sqrt{5}+i(-1+\sqrt{5})$ & $-2$ \\
6 & $2$ & $-3+\sqrt{5}+i(-1+\sqrt{5})$ & $-3+\sqrt{5}+i(1-\sqrt{5})$ & $-2$ \\
\hline
\end{longtable}
\end{center}
\end{art}

\begin{art}
The group of 60 was obtained in the $A$-form by Gordan in his paper.
The passage from the $\Theta$-form to the $A$-form is made as follows: let
$X$, $Y$, $Z$ be the coordinates of a point when the axes are as in the
$\Theta$-form, $X_1$, $Y_1$, $Z_1$ the coordinates of the same point when the axes
are as in the $A$-form: we may write
\[
X : Y : Z = bX_1 - aZ_1 : Y_1 : aZ_1 + bX_1,
\]
where
\[
a, b = \sqrt{\frac{\sqrt{5}-\sqrt{5}}{10}},\ \sqrt{\frac{\sqrt{5}+\sqrt{5}}{10}}.
\]
\end{art}

\begin{art}
The results are exhibited in terms of $\epsilon$, an imaginary fifth root of
unity: taking $\epsilon = \cos 72^\circ + i\sin 72^\circ$, we have
\[
\frac{\sqrt{5}-1}{4} + i\frac{\sqrt{10+2\sqrt{5}}}{4}, \quad
\frac{\sqrt{5}+1}{4} + i\frac{\sqrt{10-2\sqrt{5}}}{4},
\]
where the upper signs belong to $\epsilon$, $\epsilon^2$ and the lower to $\epsilon^3$, $\epsilon^4$.

It may be remarked that
\[
a = \sqrt{\frac{\sqrt{5}+1}{2\sqrt{5}}}, \quad b = \sqrt{\frac{\sqrt{5}-1}{2\sqrt{5}}}, \quad
\frac{b}{a} = \frac{\sqrt{5}+1}{\sqrt{5}-1}.
\]

For instance, we have in the $\Theta$-group $(A, B, C, D) = (-1, 0, 0, 1)$;
$ab \cdot cd$: and thence in the $A$-group $A_1, B_1, C_1, D_1 = (-2b, 2a, 2a, 2b)$;
$ab \cdot cd$: or say this is $(-1, \frac{b}{a}, \frac{b}{a}, 1) = (-1, \epsilon + \epsilon^4, \epsilon + \epsilon^4, 1)$; which in the Table
is given as $(-\epsilon^2, \epsilon^2 + \epsilon^3, \epsilon^2 + \epsilon^3, \epsilon^4)$; $ab \cdot cd$.

By effecting the passage to the $A$-group in this manner, we of course
obtain the proper substitution corresponding to each transformation: but
I found it easier starting from two transformations and the corresponding
substitutions, to obtain thence by successive compositions the entire group.
\end{art}

\begin{art}
\textbf{Homographic Transformations. The group of 60.}

Pole at~A.

\begin{center}
\begin{longtable}{|c|c|c|c|c|}
\hline
No. & $(As+B) \div (Cs+D)$ & & & \\
\hline
1 & $s$ & 1 & & \\
2 & $-s$ & $-1$ & 1 & \\
3 & $\frac{1}{s}$ & 1 & & \\
4 & $\frac{-1}{s}$ & $-1$ & & \\
5--16 & 2 & \ldots & & $-2$ \\
17 & $-1$ & & 1 & \\
18 & $i$ & & & \\
19 & $-i$ & & 1 & \\
20 & $-i$ & $-i$ & 1 & $-1$ \\
21 & $i$ & & 1 & \\
22 & $i$ & & $-i$ & \\
23 & $-1$ & $-i$ & 1 & \\
24 & $-i$ & & 1 & \\
25 & $-1-\sqrt{5}+i(3+\sqrt{5})$ & 2 & $-2$ & \\
\hline
\end{longtable}
\end{center}
\end{art}

\end{document}
